\chapter[La Física del Alma]{La Física del Alma\\ \large Y la Ley Universal del Eco.}

\elegantimage{00_introduccion/web/assets/intro.png}

\lettrine[lines=3, nindent=0.5em, findent=0.2em]{S}{}i te detienes frente a la gran pagoda de Linh Ung, en Da Nang, bajo la mirada serena de la imponente Dama Buda, encontrarás en sus muros una serie de ilustraciones cautivadoras. Son los Tranh Nhan Qua, pósters tradicionales que ilustran la inquebrantable Ley de Causa y Efecto: el Karma.

Al observar estos dibujos, es fácil pensar que estamos ante un concepto exclusivamente oriental. Sin embargo, si rascamos la superficie, descubriremos que estos murales describen una intuición que ha cruzado todas las fronteras, todas las religiones y, sorprendentemente, incluso las fronteras de la ciencia moderna.

\section{El Hilo Conductor de las Tradiciones}
La idea de que el universo lleva una contabilidad moral exacta no pertenece a una sola cultura; es el consenso espiritual de la humanidad. Aunque el término "Karma" nació en oriente, su mecanismo es universal.

\textit{Jesús y la Regla de Oro: "Así que en todo, tratad a los demás tal y como queréis que ellos os traten a vosotros" (Mateo 7:12). Generad la causa exacta que deseáis como efecto.}

\elegantimage{00_introduccion/web/assets/seed_harvest.png}

\textit{La medida: "Porque con el juicio con que juzguéis, seréis juzgados, y con la medida con que midáis, os será medido" (Mateo 7:2). Es una declaración directa de reciprocidad universal.}

Años más tarde, San Pablo cristalizaría esta idea en una máxima impecable: "Todo lo que el hombre sembrare, eso también segará" (Gálatas 6:7). Esta siembra y cosecha resuena en el Middah keneged middah del judaísmo dándole un eco eterno a nuestras acciones.

\section{La Tercera Ley de Newton del Alma}
Para comprender el Karma en el siglo XXI, quizás debamos mirar los libros de física. El universo físico detesta el desequilibrio, y la ciencia nos demuestra que toda acción genera una reacción.

En la mecánica clásica, si un objeto A ejerce fuerza sobre B, este devuelve una fuerza opuesta de igual magnitud:

\begin{center}
$\vec{F}_{A} = -\vec{F}_{B}$
\end{center}

El Karma es, esencialmente, la Tercera Ley de Newton aplicada a la conciencia. Si el universo material exige un equilibrio perfecto de fuerzas, tiene sentido lógico que el tejido de la energía espiritual opere bajo una homeostasis similar.

\elegantimage{00_introduccion/web/assets/newton_soul.png}

\section{El Entrelazamiento y la Memoria}
Incluso la física cuántica nos habla de que la información nunca se destruye. A través del entrelazamiento cuántico, una "perturbación" en un lado de la red afecta a la totalidad del sistema.

Un gesto de compasión o un estallido de ira... todo es información que se inyecta en el sistema. El universo registra, procesa y, eventualmente, devuelve esa energía para restaurar el equilibrio. No hay nada muerto, todo late e influye.

\elegantimage{00_introduccion/web/assets/quantum_entanglement.png}

\section{El Propósito de este Libro}
Este libro es un viaje a través de esa ley universal, utilizando como mapa las vívidas ilustraciones de los muros de Da Nang. No estamos aquí para juzgar, sino para observar la física del alma en acción.

Cuando comprendemos que cada pensamiento, palabra y obra es un ladrillo en la construcción de nuestro destino, dejamos de ser víctimas del azar para convertirnos en los arquitectos de nuestra propia eternidad.

\vspace{1em}
\begin{center}
    \textbf{\Large Toda acción genera su propio eco.}
\end{center}
\vspace{1em}

